\documentclass[UTF8]{ctexart}
\usepackage{xeCJK}
\usepackage{geometry}
\geometry{left=2cm,right=2cm,top=2.5cm,bottom=2.5cm}
\setlength{\headheight}{12.7pt}

\usepackage{titlesec}
\titleformat{\section}
  {\normalfont\large\bfseries}
  {\thesection}
  {1em}
  {}
\titlespacing*{\section}
  {0pt}
  {3.5ex plus 1ex minus .2ex}
  {2.3ex plus .2ex}

\usepackage{amsmath}
\usepackage{amssymb}
\usepackage{amsthm}
\usepackage{enumitem}
\setlist[enumerate]{itemsep=0pt,topsep=0pt,parsep=0pt,leftmargin=0pt,itemindent=2.5em}
\setlist[itemize]{itemsep=0pt,topsep=0pt,parsep=0pt,leftmargin=0pt,itemindent=2.5em}
\usepackage{fancyhdr}
\pagestyle{fancy}
\fancyhf{}
\usepackage{graphicx}
\usepackage{hyperref}
\usepackage{indentfirst}
\usepackage{lmodern}


\begin{document}
\setlength{\abovedisplayskip}{1pt}
\setlength{\belowdisplayskip}{1pt}

\section{半群}

\hypertarget{sec:定义1.1}{}
\noindent\textbf{定义1.1 (半群)}

给定集合$G$和$G$上的二元运算$\cdot$\,, 如果
\[
    \forall x,y,z\in G:\,x(yz)=(xy)z,
\]
就称$(G,\cdot)$是一个\textbf{半群}, 有时简写为$G$. 进一步,
\begin{enumerate}
    \item 如果$\forall x,y\in G:\,xy=yx$, 就称$(G,\cdot)$是一个\textbf{交换半群},
    \item 如果存在$e\in G$使得$\forall x\in G:\,xe=ex=x$,
    就称$(G,e,\cdot)$是一个\textbf{幺半群}, 称$e$为其\textbf{单位元}.
\end{enumerate}

\vspace{10pt}

\hypertarget{sec:定理1.1}{}
\noindent\textbf{定理1.1}

幺半群的单位元唯一, 即如果$(G,e,\cdot)$和$(G,e',\cdot)$都是幺半群, 那么$e=e'$.

\noindent{\kaishu 证明.}

如果$e,e'$都是单位元, 那么$e=ee'=e'$. \hfill $\qedsymbol$

\vspace{10pt}

幺半群对应于
单点\href{../../../02 范畴论/01 范畴、函子与自然变换/01 范畴#nameddest=sec:定义1.1}
{范畴}. 一方面, 任给幺半群$(G,e,\cdot)$,
\begin{enumerate}
    \item 定义$O=1$, $M=G$, $i(0)=e$,
    \item 对任意的$x\in G$, 定义$s(f)=t(f)=0$,
    \item 对任意的$x,y\in G$, 定义$x\circ y=xy$,
\end{enumerate}
则$\mathbf{C}=(O,M,s,t,i,\circ)$是一个单点范畴.

另一方面, 任给范畴$\mathbf{C}$和对象$A$,
$(\,\mathbf{C}(A,A),I_A,\circ\,)$是一个幺半群.

\vspace{10pt}

半群的子集之间也可以做运算成为半群, 即收集它们所有元素的积.

\vspace{10pt}

\hypertarget{sec:定义1.2}{}
\noindent\textbf{定义1.2}

给定半群$(G,\cdot)$, 对任意的$A,B\subset G$定义
\[
    A\ast B=\{ab\mid a\in A,b\in B\},
\]
则$(\,\mathcal{P}(G),\ast)$也是半群.

进一步, 如果$(G,e,\circ)$是幺半群, 那么这里$(\,\mathcal{P}(G),\{e\},\ast)$也是幺半群.

\vspace{10pt}

在子集之间的运算中, 通常把单点集$\{x\}$简写为$x$, 省略$\ast$符号.





\newpage

\section{群}

\hypertarget{sec:定义2.1}{}
\noindent\textbf{定义2.1 (群)}

给定幺半群$(G,e,\cdot)$, 如果任意的$x\in G$都有逆元$y\in G$, 即
\[
    xy=yx=e,
\]
那么称$(G,e,\cdot)$是一个\textbf{群}, 显然逆元是唯一的. 交换的群也叫\textbf{Abel群}.

\vspace{10pt}

群有以下等价的定义方式.

\vspace{10pt}

\hypertarget{sec:定理2.2}{}
\noindent\textbf{定理2.2}

任给半群$G$, 它能构成群当且仅当
\begin{enumerate}
    \item 存在左单位元$e$, 使得对任意的$x$有$ex=x$,
    \item 对任意的$x$, 存在左逆元$y$, 使得$yx=e$.
\end{enumerate}

\noindent{\kaishu 证明.}

\noindent{\kaishu 必要性.} 显然.

\noindent{\kaishu 充分性.}

首先可证左逆元就是逆元: 对任意的$x$, 取它的左逆元$y$, 再取$y$的左逆元$z$, 则
\[
    xy=(zy)xy=z(yx)y=zy=e,
\]
其次可证左单位元就是单位元: 对任意的$x$,
\begin{equation*}
    xe=x(yx)=(xy)x=ex=x.
\tag*{\qedsymbol}\end{equation*}

\vspace{10pt}

\hypertarget{sec:定理2.3}{}
\noindent\textbf{定理2.3}

任给半群$G$, 它能构成群当且仅当对任意的$a,b$, 存在$x,y$使得$ax=ya=b$.

\noindent{\kaishu 证明.}

\noindent{\kaishu 必要性.} 显然取$x=a^{-1}b,\,y=ba^{-1}$即可.

\noindent{\kaishu 充分性.}

任取$a$, 并取$e$使得$ea=a$. 此时可证$e$是左单位元: 对任意的$b$, 取$x$使得$ax=b$, 则
\[
    eb=e(ax)=(ea)x=ax=b,
\]
其次, 显然任意元素都有左逆元. 所以由上述定理得$G$是群. \hfill $\qedsymbol$

\vspace{10pt}

\hypertarget{sec:定理2.4}{}
\noindent\textbf{定理2.4}

任给有限半群$G$, 它能构成群当且仅当两侧消去律都成立.

\noindent{\kaishu 证明.}

\noindent{\kaishu 必要性.} 显然.

\noindent{\kaishu 充分性.}

如果左消去律成立, 那么对任意的$a,b$, 映射
\[
    G\to G,\quad x\mapsto ax
\]
是单射, 从而一定是满射, 即存在$x$使得$ax=b$.

同理如果右消去律成立, 那么对任意的$a,b$, 存在$y$使得$ya=b$.
所以由上述定理得$G$是群. \hfill $\qedsymbol$

\end{document}